% ============================================
% Results
% ============================================
\section{Results}
\label{sec:results}

% --------------------------------------------


\subsection{Community-level selection is prevalent in microbial coalescence}
\label{sec:results:dominance_prevalent}

Our central question is whether selection acts primarily on whole communities as cohesive units during community coalescence. To address this, we asked how often coalescence outcomes are represented primarily by one parental community versus blending of both or forming a novel community. Consider a coalescence event where communities A and B merge to produce community C (\figref{fig:fig1}a). \rev{We represent each community by a community composition vector and compute the cosine similarity between the coalesced community C and each parental community (\figref{fig:fig1}b):}
\begin{equation}
\rev{\text{Sim}(C, A) = \frac{\vec{x}_C \cdot \vec{x}_A}{\|\vec{x}_C\|_2 \|\vec{x}_A\|_2}, \quad \text{Sim}(C, B) = \frac{\vec{x}_C \cdot \vec{x}_B}{\|\vec{x}_C\|_2 \|\vec{x}_B\|_2}}
\label{eq:similarity}
\end{equation}
\rev{where $\vec{x}_A$, $\vec{x}_B$, and $\vec{x}_C$ denote community composition \revsecond{(relative-abundance)} vectors for the two parental communities and the coalesced community, and $\|\cdot\|_2$ denotes the Euclidean norm.} These two similarity scores yield a two-dimensional similarity map in which the location of C reflects how strongly it resembles each parental community. We partition this map into three coalescence outcome classes (see Methods): Dominance, where C closely resembles one parental community but not the other; Mixture, where C similarly resembles both parental communities; and Restructuring, where C diverges from both parental communities into a novel configuration. \rev{Restructuring denotes low parental retention: the coalesced community is not well explained by the parental communities or their combination and may reflect a new stable composition produced by previously unseen cross-community species interactions.} \revsecond{Throughout, Dominance is the outcome-level signature that quantifies origin-correlated persistence from compositional structure, and community-level selection is our description of that persistence rather than a claim about a specific underlying mechanism.}

To quantify the occurrence of outcome types experimentally, we assembled random \rev{\textit{in vitro}} communities and subjected them to pairwise coalescence (\figref{fig:fig1}c). We first curated a strain library of 54 bacterial isolates collected from diverse environments (soil, tree surface, and flower stamen). \revsecond{Analysis of their 252--253-bp 16S rRNA amplicon sequences identified 50 unique ASVs spanning 21 families across four phyla: Proteobacteria, Firmicutes, Bacteroidota and Actinobacteriota} (Extended Data Fig.~1; Supplementary Fig.~1). From this library, we assembled 30 parental communities with varying initial richness (6, 12, or 24 species; see Methods) and stabilized them for 7 days under daily growth--dilution cycles ($\times$30 every 24~h) (\figref{fig:fig1}c). Our initial experiments were done in Base medium (1~g~L$^{-1}$ yeast extract, 1~g~L$^{-1}$ soytone, 10~mM sodium phosphate, 0.1~mM CaCl$_2$, 2~mM MgCl$_2$, 4~mg~L$^{-1}$ NiSO$_4$, 50~mg~L$^{-1}$ MnCl$_2$, 5~g~L$^{-1}$ glucose and 4~g~L$^{-1}$ urea; pH~6.5), a buffered complex medium that we have previously determined has moderate interaction strength and supports high species coexistence (mean species survival ratio = $74 \pm 2$\%). We then performed 83 pairwise coalescence events by mixing stabilized parental communities in equal volumes and restabilizing for an additional 7 days. Community compositions of parental and coalesced communities were measured at the end of stabilization by 16S rRNA amplicon sequencing (Amplicon Sequence Variants, ASVs), and we also recorded the optical density (OD$_{600}$) and pH of the communities to contextualize assembly and post-coalescence dynamics.

Representative time series illustrate the spectrum of post-coalescence dynamics (\figref{fig:fig1}d). In outcomes classified as Dominance, one parental lineage rapidly displaces the other after mixing and the coalesced community converges to a composition closely matching that parental community, while largely preserving its internal relative-abundance structure over serial dilution cycles. In the Restructuring case, the merged community converges toward a novel state distinct from both parental communities. Among the three representative time-series trajectories, we did not observe Mixture cases in which both parental communities remained comparably represented after coalescence (see Supplementary Fig.~2).

Projecting all outcomes into the similarity space based on stabilized compositions (\figref{fig:fig1}e) revealed that Dominance is the most frequent empirical outcome. Of the 83 coalescence events, 54 were classified as Dominance, 26 as Restructuring, and 3 as Mixture (\figref{fig:fig1}e). \rev{Alternative metric analysis revealed that Dominance remained the most frequent outcome under abundance-weighted compositional metrics, including Euclidean distance and Bray--Curtis dissimilarity, but not under Jaccard index or Jensen--Shannon divergence, which emphasize species-identity retention and full-distribution divergence, respectively (Extended Data Fig.~2).} \rev{To test whether Base-medium Dominance could arise from a compositional artifact of parental abundance structure, we compared observed outcomes with matched additive-null expectations constructed from the same parental-community pairs. Observed outcomes were more strongly biased toward one parental community than their matched additive-null expectations (one-sided paired Wilcoxon test, $p = 5.9 \times 10^{-14}$; Extended Data Fig.~3), and additional abundance-skew null models also produced weaker parental bias than the data (Supplementary Note~1). \revsecond{Together, these null-model controls indicate that the observed bias toward Dominance is stronger than what parental abundance structure alone would produce.}}


% --------------------------------------------
\subsection{Theoretical model with random interactions reproduces community-level selection}
\label{sec:results:glv_model}
To gain insight into why Dominance is the prevalent outcome, we \revsecond{studied} a generalized Lotka--Volterra (gLV) model \citep{May1972, Grilli2017} that mirrors the experimental protocol (\figref{fig:fig2}a). In this classic model, species grow logistically and compete pairwise:
\begin{equation}
\frac{dn_i}{dt} = n_i \left( r_i - \sum_j \alpha_{ij} n_j \right)
\label{eq:glv}
\end{equation}
where $n_i$ is abundance, $r_i$ is growth rate, and $\alpha_{ij}$ is the competition coefficient between species $i$ and $j$ (with self-interaction $\alpha_{ii} = 1$). We fixed $r_i = 1$ and drew off-diagonal competition coefficients from a uniform distribution $\mathbb{U}(0, 2\mu)$, \rev{so increasing $\mu$ expands the sampled coefficient range and increases both the mean and spread of competitive effects} \citep{Hu2022,Hu2025}. From a pool of 54 species, \revsecond{each replicate generated four non-overlapping 12-species parental communities and coalesced all six pairs, reflecting the experimental workflow of parental-community assembly followed by pairwise coalescence} (\figref{fig:fig2}a; see Methods for details). Post-coalescence compositions were analyzed using the same similarity metrics as in the experiments. \rev{The gLV model serves as a phenomenological framework to explore how the statistical properties of interaction coefficients shape coalescence outcomes, rather than as a mechanistic model of the specific biochemical interactions (e.g., pH modification, metabolic cross-feeding) underlying our experimental observations.}

\revsecond{We repeated this four-community procedure 200 times at} \rev{an interaction-strength parameter value of} $\mu = 0.6$, \revsecond{yielding $200 \times 6 = 1{,}200$ random coalescence events,} and analyzed the similarity of each coalescence outcome to the two parental communities. At this \rev{parameter value}, the random-interaction model quantitatively reproduces the high frequency of Dominance observed in the experiments. Outcomes concentrate near the axes rather than the diagonal, indicating asymmetric dominance by one parental community (\figref{fig:fig2}b). Quantitatively, Dominance accounts for 61\% of all outcomes, far more frequent than Restructuring (26\%) or Mixture (13\%; \figref{fig:fig2}b, right). This observation suggests that, within this competitive gLV framework, interspecies competitive interactions are sufficient to reproduce the high frequency of Dominance observed in coalescence.

To understand how Dominance emerges in the model and whether it reflects community-level selection, we examined the role of the assembly process. During assembly, competitive exclusion filters out species with strong competitive interactions, leaving communities with reduced \rev{interaction strength} compared to the initial pool (paired $t$-test, $P <0.001$; \figref{fig:fig2}c; Supplementary Fig.~3). Because surviving species compete weakly with each other but face stronger competition from outsiders, their fates become coupled during coalescence: they tend to persist or go extinct together. \rev{We quantified this effect using pairwise selection correlation, which tests across coalescence events whether species pairs from the same parental community share survival outcomes (both persist or both go extinct) more strongly than species pairs from different parental communities (Supplementary Note~7; \figref{fig:fig2}d).} Species from the same parental community showed strongly positive correlations, while cross-community pairs showed negative correlations, a pattern observed in both simulations and experiments (\figref{fig:fig2}d). \revsecond{By comparing coalescence with direct assembly of the same combined species pool, in which species are mixed without prior separate assembly, we found that coalescence produced substantially more Dominance in both simulations and the corresponding experimental conditions (59\% versus 20\% at $\mu = 0.6$ in simulations; Supplementary Note~2; Extended Data Fig.~4; Supplementary Fig.~23). This excess therefore reflects prior assembly history. Together, these results support the interpretation that species within an assembled parental community have coupled fates during coalescence, the origin-correlated persistence we describe as community-level selection.}


% --------------------------------------------
\subsection{Interaction strength controls coalescence outcome type and the degree of community-level selection}
\label{sec:results:interaction_strength}

Having established that random interspecies interactions can recapitulate Dominance with correlated pairwise selection, we next investigated how the interaction-strength parameter alters coalescence outcomes. We simulated 1,200 coalescence events at each of three representative \rev{values of the interaction-strength parameter} ($\mu = 0.3, 0.6, 0.8$) and mapped outcomes into the similarity space (\figref{fig:fig3}a). At weak interactions ($\mu = 0.3$), outcomes clustered in the interior of the map, indicating Mixture; as $\mu$ increased, outcomes shifted toward the axes, indicating Dominance. We quantified this shift using the Parental Dominance Index (PDI), which captures the relative contribution of each parental community (0: parental community B, 0.5: equal, 1: parental community A; see Methods). PDI distributions shifted from unimodal near 0.5 at $\mu = 0.3$ to bimodal at higher $\mu$ (\figref{fig:fig3}a), reflecting a transition in which Dominance becomes increasingly frequent, indicating community-level selection.

Across \rev{interaction-strength parameter values} from $\mu = 0$ to $1.2$, we observed a systematic shift in coalescence outcomes (\figref{fig:fig3}b), with Mixture predominating at weak interactions and Dominance at strong interactions. Restructuring emerged at moderate to high interaction strengths. \rev{These patterns were robust to variation in growth rates, carrying capacities, interaction-coefficient distributions, parental-community species richness, and model extensions allowing beneficial interactions or varying reciprocal pair-coupling structure (Supplementary Figs.~4--6 and 39--41; Extended Data Fig.~5). Further analysis separating the mean and variance of the interaction-coefficient distribution showed that both the average inhibitory effect and coefficient heterogeneity contribute to the Dominance regime (Supplementary Fig.~42).} \revsecond{The frequency of Dominance remained relatively stable across parental communities ranging from 4 to 48 species, spanning our experimental species pool range.} To determine whether this transition from Mixture to Dominance corresponds to the degree of co-selection, we examined pairwise selection correlation across \rev{interaction-strength parameter values}. \revsecond{Within-community pairwise selection correlation emerged only at high interaction strengths, indicating that species fates become coupled only when interactions are sufficiently strong} (Extended Data Fig.~6). These findings indicate that interaction strength simultaneously determines both coalescence outcome type and the level at which selection operates.


% --------------------------------------------
\subsection{Nutrient-dependent invasion resistance and environmental feedbacks shift coalescence outcomes}
\label{sec:results:nutrient_experiment}

\rev{Following prior work showing that nutrient concentration can alter microbial competition and that pH-mediated environmental feedbacks can drive bacterial interactions \citep{Hu2022,Ratzke2018,Ratzke2020,Hu2025}, we conducted additional coalescence experiments by removing glucose and urea from the Base medium or increasing their concentrations (Methods). This yielded two additional media conditions (\figref{fig:fig4}a): Nutr$-$ (no added glucose/urea) and Nutr$+$ (high supplementation). To approximate \revsecond{changes in net effective interaction strength} across nutrient conditions, we measured invasion resistance using pairwise invasion assays among the 12 most abundant isolates (95:5 initial frequency; Methods, Supplementary Figs.~7--9). \revsecond{Failed-invasion frequency increased with nutrient supply} (Nutr$-$: $2 \pm 1\%$; Base: $33 \pm 4\%$; Nutr$+$: $48 \pm 4\%$; \figref{fig:fig4}b), \revsecond{indicating greater invasion resistance and potentially stronger competition or other antagonistic effects, including those mediated by environmental feedbacks, at higher nutrient levels \citep{Ratzke2018,Ratzke2020}}.}

\rev{We then performed coalescence experiments in Nutr$-$ and Nutr$+$ media using the same parental-community assembly scheme to examine whether this nutrient-dependent change in invasion resistance is accompanied by the \revsecond{expected shift in coalescence outcomes, from Mixture toward Dominance as effective interactions strengthen (\figref{fig:fig3})}.} The distribution of outcomes in similarity space shifted systematically with nutrient concentration (\figref{fig:fig4}c), which we quantified by the fraction of each outcome type (\figref{fig:fig4}d). In Nutr$-$ ($n = 90$), where interactions are weakest, Mixture was the most frequent outcome (53\%), and Dominance was substantially reduced to 39\% compared to the Base medium (Dominance: 65\%, Mixture: 4\%). In Nutr$+$ ($n = 90$), Dominance further increased to 76\%. \rev{The distributions of PDI, asymmetry, and retention magnitude show the same nutrient-dependent shift \revsecond{before classifying outcomes categorically} (Supplementary Figs.~10--12).} Pairwise selection correlation analysis confirmed that this shift corresponds to a transition in the level of selection (Extended Data Fig.~7). In Base and Nutr$+$ media, within-community species pairs showed significantly higher selection correlation than cross-community pairs ($P < 0.001$), consistent with community-level selection. In Nutr$-$, no such correlation was observed, consistent with weaker origin-correlated persistence and more species-level dynamics. \revsecond{These results match the qualitative theoretical prediction and extend it to the level of selection: weak effective interactions yield Mixture with uncorrelated species fates, whereas strong effective interactions yield Dominance with community-level selection.}

\revsecond{The nutrient-dependent shift could alternatively reflect changes in parental-community properties unrelated to interaction strength.} \rev{We \revsecond{therefore} tested whether nutrient-dependent changes in parental-community properties, including parental-community biomass heterogeneity and initial richness effects, could explain this trend; however, Dominance outcomes were not biased toward the higher-biomass parental community, and these factors did not account for the nutrient-dependent increase in Dominance or correlated species fates (Supplementary Note~5 and Supplementary Figs.~13--15, 45).}


% --------------------------------------------
\subsection{Predictability of Dominance direction reveals two mechanistic regimes}
\label{sec:results:predictability}

We observed Dominance in both Base and Nutr$+$ media, consistent with community-level selection as confirmed by pairwise selection correlation (Extended Data Fig.~7). We next asked whether we could predict which parental community would win during coalescence. Prior work has suggested that species-level competitive outcomes may correlate with Dominance direction \citep{Rillig2015, Castledine2020}. Inspired by these observations, we tested whether pairwise competition between dominant species \citep{Hu2022, Ratzke2020} predicts which community wins (\figref{fig:fig5}a).

The relative abundance of dominant species in stabilized parental communities increased with nutrient concentration ($44 \pm 2\%$ in Nutr$-$, $51 \pm 5\%$ in Base, $67 \pm 4\%$ in Nutr$+$; mean $\pm$ s.e.m., \figref{fig:fig5}b), suggesting that dominant species exert greater influence at higher nutrient concentrations. We therefore tested whether pairwise competition between dominant species predicts the coalescence outcome by analyzing the linear relationship between dominant-species competitive success (from invasion assays) and the PDI of coalescence outcomes (\figref{fig:fig5}c). Predictive power varied markedly across media: in Nutr$-$, where Mixture predominates and pairwise selection correlation is weak, pairwise competition showed no predictive power ($R^2 = 0.00$); in Base medium, predictive power was weak ($R^2 = 0.11$), consistent with collective multi-species dynamics shaping outcomes; in Nutr$+$, pairwise competition was more predictive ($R^2 = 0.49$), indicating that dominant species contribute substantially to determining which community wins. These results suggest that community-level selection spans a mechanistic continuum: at one end, an emergent regime (Base medium) where multi-species dynamics collectively shape the outcome and no single species trait is predictive, and at the other, a top-down regime (Nutr$+$ medium) where a few dominant taxa determine the winner. \rev{The Nutr$+$ predictability trend remained significant after recalculating PDI with dominant species excluded (Spearman rank-correlation tests: $p = 8.2 \times 10^{-4}$ after removing the coalesced-community dominant species and $p = 1.3 \times 10^{-3}$ after removing the two parental-community dominant species; Supplementary Fig.~16), indicating that \revsecond{the predictive power of dominant-species interaction outcomes is not simply a result of including dominant species in the PDI calculation}.}

\rev{We next asked whether environmental modification through pH could contribute to the top-down regime.} In our system, dominant species are often strong pH modifiers that either acidify or alkalinize the medium (Supplementary Fig.~17). In both Base and Nutr$+$ media, when these pH-modifying species become dominant within a community, \rev{their identities predict} the community's overall pH. Communities dominated by acidifiers become acidic, while those dominated by alkalizers become alkaline (Supplementary Fig.~18). \rev{In coalescence events between acidic (pH $<$ 6.5) and alkaline (pH $>$ 7.5) parental communities, the acidic community dominated in only 56\% of cases in Base medium ($n = 41$) and in 91\% of cases in Nutr$+$ medium ($n = 32$) (Extended Data Fig.~8). The difference between Base and Nutr$+$ proportions was significant (two-sided Fisher's exact test, $p = 0.0015$), consistent with the hypothesis that high nutrient levels amplify metabolic activity, thereby intensifying pH modification and excluding pH-sensitive species \citep{Ratzke2018, Ratzke2020}.} \rev{Together, the dominant-species and strict pH-contrast analyses suggest that Nutr$+$ Dominance may reflect a top-down regime, in which dominant taxa and associated environmental feedbacks help determine winner identity. However, this pH-based analysis addresses winner direction in high-contrast acidic--alkaline pairings and does not by itself account for the full Dominance pattern across all coalescence events. Dominance in Base medium is less directly explained by the measured dominant-species and pH effects tested here.}


% --------------------------------------------
\subsection{\texorpdfstring{\rev{Nutrient-dependent coalescence outcomes in natural sample-derived communities}}{Nutrient-dependent coalescence outcomes in natural sample-derived communities}}
\label{sec:results:natural_communities}

The synthetic communities described above were assembled from individually isolated bacterial strains, allowing precise control over initial composition but raising the question of whether these patterns also appear in taxonomically richer communities derived from natural environmental samples. \rev{To address this, we performed coalescence experiments using natural sample-derived communities established through laboratory stabilization from environmental samples.}

We collected six environmental samples from diverse microhabitats (soil, compost, and decomposing organic matter) and established bacterial communities through seven rounds of serial growth-dilution in laboratory media across all three nutrient conditions (Nutr$-$, Base, and Nutr$+$; \figref{fig:fig6}a). After stabilization, these natural sample-derived communities exhibited higher ASV richness than synthetic communities (mean of $13.7 \pm 7.2$ ASVs above the 0.1\% threshold, compared to $9.8 \pm 4.8$ in synthetic communities; Supplementary Figs.~19--21) and \rev{retained source-specific taxonomic structure at both ASV and genus levels (Supplementary Fig.~22).}

We performed 15 pairwise coalescence events with two biological replicates each ($n = 30$ per condition) across all three nutrient conditions after stabilization and analyzed outcomes using the same similarity framework applied to synthetic communities. Natural sample-derived communities showed similar nutrient-dependent patterns to synthetic communities: outcomes were more centered in Nutr$-$ and shifted toward the axes in Base and Nutr$+$ (\figref{fig:fig6}b,c). The fraction of Dominance outcomes increased with nutrient concentration (37\% in Nutr$-$, 70\% in Base, 77\% in Nutr$+$), mirroring the nutrient-dependent trend observed in synthetic communities. \rev{These communities showed higher Restructuring fractions, possibly reflecting greater taxonomic diversity and more complex interaction networks. Because these communities were laboratory-stabilized before coalescence, pre-selection in defined media could have contributed to their similarity to synthetic-community outcomes. However, their higher richness and retained source-specific taxonomic structure indicate that stabilization did not produce complete ASV- or genus-level taxonomic collapse (Supplementary Figs.~19--22 and Supplementary Note~6). Overall, these results suggest that the nutrient-dependent increase in Dominance observed in synthetic consortia is also present in taxonomically richer natural sample-derived communities.}
